\selectlanguage{english}
\noindent\textbf{Chapter VII, manuscript pages 1--57: the shapes of paths.}

The opening of Chapter VII is a restart. Grothendieck has found the earlier path formalism too cumbersome to serve as a foundation for localization, and he now separates a path from the finite combinatorial object on which it is parametrized. The distinction gives these pages their unity. First one understands intervals independently of any coefficient category; next one regards every operation on paths as a change of interval; finally one assembles all types and makes their concatenation strictly functorial. The later theory of fractions is a test of the construction, not yet a theorem established in these pages. The commentary on localization and path completion in Chapter~VI, especially at manuscript pages~64--70, supplies the background; the present chapter changes the finite apparatus on which such arguments are to rest.

An interval of length $n$ is the free category on an oriented path with vertices $0,\ldots,n$. Its signature is a word $\tau=(\varepsilon_1,\ldots,\varepsilon_n)$ in two signs. There are two orders to remember. The progression order is the linear order of the vertices; the categorical order records the existence of arrows. A path proceeds from vertex~$0$ to vertex~$n$ even when some of its arrows point backwards. Reversing this progression order changes the signature to $(-\varepsilon_n,\ldots,-\varepsilon_1)$, whereas taking the opposite category changes it to $(-\varepsilon_1,\ldots,-\varepsilon_n)$. Their distinction is already needed to understand why a span is symmetric as an interval but is not isomorphic to its opposite cospan.

The maximal constant runs of the signature are the structural chains. Their endpoints are the critical vertices; retaining just these vertices gives the reduced interval. These notions organize all full subcategories of an interval. A nonempty full subcategory has interval components, and a subset with extreme vertices $i<j$ is connected precisely when it retains every critical vertex between them. The proof is elementary and decisive: two consecutive retained vertices can be joined by an arrow exactly when no change of sign has been omitted between them. Thus reduction remembers the turns while suppressing subdivisions inside a composable run. It does not mean that arbitrary diagrams on the interval may forget the intermediate objects.

Several literal statements near the beginning need repair. Proposition~1(b) asserts that every endpoint-fixing endofunctor is the identity. On $[2]$, the function $(0,1,2)\mapsto(0,0,2)$ is a counterexample. The proof on the following pages establishes rigidity of endpoint-fixing \emph{isomorphisms}, which is sufficient for the classification of orientations. In fact, the stronger endofunctor assertion holds exactly for reduced intervals. Proposition~3's deletion criterion also needs the vertices incomparable with the deleted vertex to be separated from its incoming and outgoing links. Finally, an interval-shaped subcategory need not be full. All these readings occur in the edition and remain in the source text. The accompanying blocks give counterexamples and prove the repaired assertions, so that subsequent correct arguments do not inherit an unnecessary uncertainty.

Reduction has a useful additional interpretation. For the full inclusion $j:I_{\mathrm{red}}\hookrightarrow I$, send an interior vertex of each run to that run's categorical target or categorical source, and fix all critical vertices. The resulting maps satisfy
\[
 p_{\mathrm{target}}\dashv j\dashv p_{\mathrm{source}}.
\]
The commentary proves these adjunctions by comparing the relevant hom-sets. They give categorical deformation retractions and explain exactly why reduction is harmless for groupoid-valued paths. For a general category $X$, restriction $X^I\to X^{I_{\mathrm{red}}}$ need not be an equivalence. This is the first place where the geometry of the shape and the information in a labelled diagram must be kept distinct.

Once the shape is fixed, a path of type $\tau$ in $X$ is a functor $I_\tau\to X$, and its morphisms are natural transformations. The category is $X^{I_\tau}$; its set of objects is the represented functor $\Hom_{\Cat}(I_\tau,X)$. Yoneda then says that every operation on paths natural in the coefficient category is precomposition by a unique functor between their shapes, in the reverse direction. The proof evaluates the operation on the identity path in $I_\tau$. This explains both composition of successive arrows and insertion of identities, and supplies their action on morphisms without an additional coherence choice. The manuscript's distinction between the set and the category of paths must therefore survive the typography.

The category of types changes twice while Grothendieck tests it. Initially it permits all functors between intervals; then it retains functors preserving the two endpoints; finally it also requires their vertex functions to be nondecreasing in the progression order. Only the first version is fully faithful in $\Cat$. The second is fully faithful in the category of categories under the discrete two-point category. The final version, denoted $\uTy$, is generally not full there either. These are successive definitions, not contradictory assertions about a fixed category. The final restriction preserves directed progress through a path and excludes changes of parameter which backtrack.

For any chosen version $T$, the strict contravariant functor $\tau\mapsto X^{I_\tau}$ gives a split Grothendieck fibration over $T$. A morphism from $(\tau',c')$ to $(\tau,c)$ is a pair
\[
 \varphi:\tau'\to\tau,
 \qquad v:c'\Rightarrow c\varphi.
\]
Composition is whiskering followed by composition of transformations. This is the machinery of \citet[Expos\'e~VI, \S\S8--10]{GrothendieckSGA1}, specialized to restriction along finite shapes. Endpoint-preserving shape maps make evaluation at the two ends into functors on the total category. The commentary also determines the maximal enlargement for which natural endpoint evaluations can still exist, allowing comparison arrows at the ends. That enlargement illustrates why the subsequent requirement of strict gluing is stronger. No colimits in $X$ are needed for the fibration; extending paths in the opposite direction would require Kan extensions, and already a collapsed span asks for a pushout.

The main calculation of the second half has a particularly short modern form. Give the transition set $\Tr(\tau)=\{1,\ldots,n\}$ its order and its signs. If $f:I_\tau\to I_\upsilon$ is a morphism of the final type category, every target transition belongs to the image interval of exactly one source transition. This defines a monotone, sign-preserving map
\[
 f^*:\Tr(\upsilon)\longrightarrow\Tr(\tau).
\]
Conversely, given such a map $h$, recover the vertex function by
\[
 f(i)=\#\{j:h(j)\leq i\}.
\]
These assignments are inverse and reverse composition. Thus the type category is opposite to the category of finite signed linear orders with monotone sign-preserving maps. Its many normal forms are descriptions of this one correspondence.

For a reduced target, $f^*$ is injective, so a map is specified by a signed subsequence of the source word. In general, a map out of a fixed source is specified by nonnegative integer weights: an edge of sign $\varepsilon$ and weight $a$ becomes a run of $a$ edges of sign $\varepsilon$. Weight zero collapses that edge. A map into a fixed target is instead described by its image vertices and by the signatures of the fibers over them. The latter description must use cumulative sums for the target coordinates, and collapsed fibers are components of the graph of collapsed edges, not components of the union of their endpoint vertices with every ambient adjacency restored. The point target is treated separately. These corrections preserve the two classification results.

Concatenation glues the final vertex of one interval to the initial vertex of the next. On standard words and endpoint-preserving maps it is strictly associative and unital, and makes $\uTy$ a strict monoidal category. It cannot extend with this object operation to the category allowing \emph{all} interval functors: there $+$ and $-$ are isomorphic, whereas $++$ and $-+$ are not. The restriction on morphisms is therefore essential to the source's own construction. Through the signed-edge description, $\uTy$ is the free strict monoidal category on two counital coassociative comonoids. The one-generator background is the opposite of the augmented simplex category; see \citet[Appendix~D, Proposition~D.2]{AguiarMahajan2010}. The two-generator version follows by ordered concatenation, without a symmetry interchanging the letters.

The nerve calculation is now transparent. A chain of $r$ composable type maps splits uniquely according to the edges of its initial word. Each factor begins at a single signed edge; its later stages are successive ordered expansions of that edge, allowing complete disappearance. Consequently every $N_r(\uTy)$ is a free monoid. The face deleting the initial object can send a generator to several generators, or to the empty word, so this assertion does not identify the nerve with the free simplicial monoid on its degreewise generators. The commentary gives every face and degeneracy by the corresponding operation on ordered forests and resolves the final support-indexing slip.

The later comparison can be made at the level of the indexing category itself. Marking the negative transitions identifies $\uTy$ with the zigzag type category in \citet[Definition~3.5]{MazelGee2018Hammocks}. The closing status block distinguishes this finite identification from the additional construction of localization mapping spaces, including the recent comparison in \citet[Theorem~1.9 and Corollary~1.11]{ArakawaCnossen2026}. The development blocks prove explicit consequences left unworked here: the reduction adjunctions and sharp rigidity criterion, the endpoint evaluation criterion, and a criterion with a counting formula for sections of contractions. Their markers record our proofs and limited literature searches, not a claim of priority. The unresolved task at manuscript page~57 is the promised construction of fractions with its precise morphisms; the finite combinatorics considered in these pages has been settled in the commentary.
